
\begin{figure}[t]
     \centering
     \includegraphics[width=0.4\textwidth]{img/mind_tradeoff.png}
     \caption{MIND Experiment Computation-Performance Tradeoff. Better Tradeoff towards Top Right. We expect full-scale NeuralUCB and NeuralTS to perform similarly to our methods, but require orders of magnitude more compute. }
     \label{fig:mind_tradeoff}
     \vspace{-0.15in}
 \end{figure}

Recommender systems (RS), paramount in personalizing digital content, critically influence the quality of information accessed via the Internet. Traditionally, these systems have employed supervised learning algorithms, such as Collaborative Filtering \cite{resnick1994grouplens}, which have greatly benefited from advances in deep learning. These algorithms analyze vast quantities of data to discern user preferences; however, they are not designed to strategically probe in order to more quickly learn about user interests. Instead, they learn passively from collected data.

Current research \cite{jiang2019degenerate, salamo2005reducing, gravino2019towards} reveals that deep-learning-driven RS tend to quickly confine their focus to a limited set of suboptimal topics, limiting their scope and hampering their learning capacity. This restrictive personalization strategy confines RS to recommend only those topics with which they have established familiarity, thus failing to discover and learn users' other potential interests. The ability of an RS to identify and learn about user's unexplored interests is a significant determinant of its long-term performance.

The concept of exploration in RS draws heavily from bandit learning \cite{slivkins2019introduction}. In the bandit learning formulation of RS, the system functions as an agent, each user represents a unique context, and each recommendation is an action. Bandit learning algorithms, such as upper confidence bound (UCB) \cite{agrawal1995sample, auer2002finite} and Thompson sampling \cite{thompson1933likelihood}, provide the groundwork for efficient exploration.  Theoretical advances \cite{agrawal2012analysis, dong2018information, dong2019performance, auer2002using, agrawal1995sample, auer2002finite} have offered great insight into these methods, and the efficacy of such methods has been demonstrated through experiments with small-scale environments.  While much of this literature has focused on small-scale environments that do not require that an agent generalizes across contexts and actions, the scale of practical RS calls for such generalization.

Advances in neural bandits offer flexible and scalable approaches to generalization that support more sample-efficient exploration algorithms.  Building on ideas developed for linear UCB and linear Thompson sampling \cite{li2010contextual, chapelle2011empirical}, these approaches offer variations of UCB and Thompson sampling that interoperate with deep learning \cite{li2010contextual, zhou2020neural, xu2021neural,abeille2017linear, zhang2021neural, lu2017ensemble,dwaracherla2019hypermodels,OsbandNeuralTestbed,QinEnsembleSampling}.  While these approaches may be sample-efficient, their computational requirements become onerous at scale.  This has limited their practical impact.

An obstacle to scaling aforementioned approaches has been in the computation required to maintain and apply epistemic uncertainty estimates.  Such estimates allow an agent to {\it know what it does not know}, which is critical for guiding exploration.  The EpiNet \cite{osband2021epistemic} offers a scalable approach to uncertainty estimation, and therefore, a path toward supporting efficient exploration in practical RS.  In particular, epinet-enhanced deep learning combined with Thompson sampling, has the potential to greatly improve RS' exploration capabilities and therefore improve personalization. To this end, we introduce Epistemic Neural Recommendation (ENR), a novel architecture customized for RS. We run a series of experiments using large-scale real-world RS datasets to empirically demonstrate that ENR outperforms state-of-the-art neural contextual bandit algorithms. ENR greatly enhances personalization, achieving an $9\%$ and $6\%$ improvement in click-through rate and user rating, respectively, across two real-world experiments. Furthermore, it attains the performance of the best baseline algorithm with at least $29\%$ fewer user interactions. Importantly, ENR accomplishes these while requiring orders of magnitude fewer computational resources than neural contextual bandit baseline algorithms, making it a considerably more scalable solution. As a spoiler, please see Figure \ref{fig:mind_tradeoff} for a computation-performance tradeoff comparison between our method and baseline methods based on one of our real-world experiments.


% ENR not only captures user preferences more accurately by improving click through rate and user rating by $18\%$ and $6\%$ in two real-world experiments, but also achieves the same performance as the best state-of-the-art algorithm with at least $20\%$ fewer user interactions. Moreover, ENR achieves this performance while requiring orders of magnitude less computational resources.
% Each of these datasets offers full ground truth user feedback for all available user-recommendation pairs and hence no counterfactual evaluation is needed to evaluate algorithm performance. 

The remainder of this paper is organized as follows. In Section 2, we formulate RS as a contextual bandit. In Section 3, we review existing online contextual bandit algorithms, current industry adoptions, and their relative merits. In Section 4, we introduce Epistemic Neural Recommendation. In Section 5, we present empirical results.  In Section 6, we summarize our contributions, benefits they afford, and potential future directions for neural contextual bandits.

